\section*{الفصل \ref{ch:b1:plant-water-minerals} --- تغذية النبات بالماء والمعادن}

\begin{solution}{exo:b1:plant-water-minerals:1}
\omterm{def:b1:plant-water-minerals:pathways}{الأبوبلاست}: الجدر والمسافات المتصلة، تُعبر دون المرور بغشاء.
\omterm{def:b1:plant-water-minerals:pathways}{والسمبلاست}: \omterm{def:b1:cell-unit-of-life:cell}{سيتوبلازم} \omterm{def:b1:cell-unit-of-life:cell}{الخلايا} الموصول، يُدخل إليه عبر غشاء واحد
ويستمر عبر \omterm{def:b1:eukaryotic-cell:plantcell}{الروابط السيتوبلازمية}. \omterm{def:b1:plant-water-minerals:pathways}{وشريط كاسباري} يسدّ \omterm{def:b1:plant-water-minerals:pathways}{الأبوبلاست} عند
\omterm{prop:b1:flowering-plant-organization:stemroot}{البشرة الداخلية}، فيُلجئ كل شيء إلى المرور بغشاء ويمنع التسرب
الراجع من \omterm{def:b1:flowering-plant-organization:tissues}{الخشب}.
\end{solution}

\begin{solution}{exo:b1:plant-water-minerals:2}
N (\omterm{def:b1:proteins:peptide}{البروتينات} والأحماض النووية)، وP (\omterm{def:b1:nucleic-acids:nucleotide}{النوكليوتيدات} \omterm{def:b1:membranes-transport:membrane}{والأغشية})، وK
(عامل \omterm{def:b1:membranes-transport:osmosis}{التناضح} والثغور)، وCa (الجدر والإشارات)، وMg (\omterm{def:b1:photosynthesis:pigments}{اليخضور})، وS
(الحمضان الأمينيان السيستين والميثيونين).
\end{solution}

\begin{solution}{exo:b1:plant-water-minerals:3}
التربة \qty{-0.2}{}، \omterm{def:b1:flowering-plant-organization:tissues}{وخشب} \omterm{def:b1:flowering-plant-organization:organs}{الجذر} \qty{-0.6}{}، \omterm{def:b1:cell-unit-of-life:cell}{وخلية} \omterm{def:b1:flowering-plant-organization:organs}{الورقة}
\qty{-1.3}{MPa}: من التربة إلى \omterm{def:b1:flowering-plant-organization:organs}{الجذر}، ومن \omterm{def:b1:flowering-plant-organization:tissues}{خشب} \omterm{def:b1:flowering-plant-organization:organs}{الجذر} صعودًا إلى
\omterm{def:b1:flowering-plant-organization:organs}{الورقة} --- دائمًا نحو القيمة الأشد سلبية.
\end{solution}

\begin{solution}{exo:b1:plant-water-minerals:4}
النتروجين متحرك في اللحاء: فالنبات يسحبه من \omterm{def:b1:flowering-plant-organization:organs}{الأوراق} القديمة ليغذي
الفتية، فتصفرّ القديمة أولًا. أما الحديد فلا يُعاد تحريكه: \omterm{def:b1:flowering-plant-organization:organs}{فالأوراق}
الفتية، وقد بُنيت بدونه، هي التي تصفرّ.
\end{solution}

\begin{solution}{exo:b1:plant-water-minerals:5}
كمون \omterm{def:b1:cell-unit-of-life:cell}{الخلية} $\Psi = -0.9 + 0.5 = \qty{-0.4}{MPa}$: فالتربة ($-0.2$) $>$ \omterm{def:b1:cell-unit-of-life:cell}{الخلية}
($-0.4$) $>$ \omterm{def:b1:flowering-plant-organization:tissues}{الخشب} ($-0.6$)، فيجري الماء من التربة $\to$ \omterm{def:b1:cell-unit-of-life:cell}{الخلية} $\to$
\omterm{def:b1:flowering-plant-organization:tissues}{الخشب}. ويتوقف الالتقاط حين تبلغ \omterm{def:b1:cell-unit-of-life:cell}{الخلية} \qty{-0.2}{MPa}، أي عند
$\Psi_p = -0.2 + 0.9 = \qty{0.7}{MPa}$.
\end{solution}

\begin{solution}{exo:b1:plant-water-minerals:6}
$\Psi = 137\ln(\text{الرطوبة النسبية})$: أي \qty{-14}{MPa} و\qty{-95}{MPa}
و\qty{-315}{MPa}. فحتى الهواء الرطب أجف بعشر مرات، بهذا المعنى، من
تربة عند نقطة الذبول: \omterm{def:b1:flowering-plant-organization:organs}{فالورقة} تفقد الماء إلى الهواء دائمًا؛ والسؤال
الوحيد هو بأي سرعة.
\end{solution}

\begin{solution}{exo:b1:plant-water-minerals:7}
$E_K = 25.7\ln(0.1/100) = \qty{-178}{mV}$: فعند \qty{-150}{mV} تكون
\omterm{def:b1:cell-unit-of-life:cell}{الخلية} أقل سلبية قليلًا مما يلزم، فتراكم بألف مثل يحتاج إلى شيء
أكثر قليلًا من \omterm{def:b1:membranes-transport:transporters}{القنوات} --- لكن نسبة 100 تكون عند التوازن: \omterm{def:b1:membranes-transport:transporters}{فالقنوات}
تفعل ذلك كله تقريبًا. أما النترات: $E = -25.7\ln(0.5/5)
= \qty{+59}{mV}$؛ \omterm{def:b1:cell-unit-of-life:cell}{والخلية} عند \qty{-150}{}، أي على بعد \qty{209}{mV}
وفي الاتجاه الخاطئ: فوجب ضخ الأنيون داخلًا، بنقل مشترك مع البروتون.
\end{solution}

\begin{solution}{exo:b1:plant-water-minerals:8}
ليلًا، وبلا نتح، يجتذب تحميل الأيونات في \omterm{def:b1:flowering-plant-organization:tissues}{الخشب} الماءَ ويدفع \omterm{prop:b1:plant-water-minerals:rootpressure}{ضغط
جذري} مقداره أعشار من الميغاباسكال العصيرَ صعودًا بضعة أمتار على
الأكثر --- يكفي لبلوغ أطراف نصال العشب وإخراج قطيرات، ولا يكفي
لبلوغ قمة شجرة. ومتى انفتحت الثغور وضع النتح \omterm{def:b1:flowering-plant-organization:tissues}{الخشب} تحت توتر وزال
\omterm{prop:b1:plant-water-minerals:rootpressure}{الضغط الجذري}.
\end{solution}

\begin{solution}{exo:b1:plant-water-minerals:9}
بلا موليبدنوم لا تعمل \omterm{prop:b1:plant-water-minerals:nitrate}{مختزلة النترات}: فعلى النترات لا يستطيع النبات
صنع \omterm{def:b1:proteins:aminoacid}{أحماض أمينية} ويظهر عوز نتروجين مع وفرة النترات. وعلى الأمونيوم
ينمو نموًا طبيعيًا، إذ يدخل الأمونيوم في \omterm{def:b1:proteins:aminoacid}{الأحماض الأمينية} بلا
اختزال. والفرق يحدد موضع العنصر في \omterm{def:b1:enzymes:enzyme}{إنزيم} واحد.
\end{solution}

\begin{solution}{exo:b1:plant-water-minerals:10}
\qty{30}{mg} من N $= \qty{2.14}{mmol}$؛ و$\times 8 = \qty{17}{mmol}$ من
الإلكترونات، أي \qty{8.6}{mmol} من مكافئات \omterm{def:b1:biosyntheses-integration:ppp}{NADPH}. \omterm{def:b1:flowering-plant-organization:organs}{والورقة}: $0.05\times
4\times 10^{-6}\times 43\,200 = \qty{8.6}{mmol}$ من $\mathrm{CO_2}$،
أي \qty{35}{mmol} من الإلكترونات: فاختزال النترات يأخذ نصف ما
يأخذه تثبيت الكربون من الإلكترونات على هذا الحساب --- وهو كسر كبير،
ولذلك تقوم به النباتات في الضوء، ولذلك تكون \omterm{def:b1:flowering-plant-organization:organs}{الأوراق} السريعة النمو
مستهلكة كبيرة للمرجع.
\end{solution}

\begin{solution}{exo:b1:plant-water-minerals:11}
\qty{2.14}{mmol} من N $= \qty{1.07}{mmol}$ من $\mathrm{N_2}$: ومقدار \omterm{def:b1:water-small-molecules:atp}{ATP}
هو $16\times 1.07 = \qty{17}{mmol}$، أي \qty{0.57}{mmol} من الغلوكوز،
أي \qty{0.10}{g}؛ زائد 8 إلكترونات لكل $\mathrm{N_2}$ (أي
\qty{0.05}{g} أخرى من الغلوكوز). \omterm{def:b1:nucleic-acids:nucleotide}{وقاعدة} \qty{6}{g/g} تعطي
\qty{0.18}{g} في اليوم: فالنتروجيناز نفسه أكثرها، والباقي بناء
\omterm{def:b1:plant-water-minerals:nodules}{العقيدات} وصيانتها، وليغهيموغلوبينها، وتنفس البكتيريا نفسها لإبقاء
الأكسجين منخفضًا.
\end{solution}

\begin{solution}{exo:b1:plant-water-minerals:12}
\omterm{def:b1:flowering-plant-organization:organs}{الجذر} نفسه يأخذ الماء والأيونات المتحركة (النترات والبوتاسيوم) التي
يجلبها ماء التربة إلى أشعاره، ولا ينفق \omterm{def:b1:water-small-molecules:atp}{ATP} إلا على \omterm{prop:b1:plant-water-minerals:uptake}{مضخة البروتونات}.
أما الفوسفات، وهو لا يأتي إليه، فيدفع لفطر سكرًا ليمد متناوله؛ وأما
النتروجين في تربة بلا نترات فيدفع لبكتيريا سكرًا لتثبته. والأجر
عُشر تركيبه الضوئي أو أكثر، ويكفّ النبات عن الدفع متى لم تكن
الخدمة لازمة: فعلى التربة الغنية بالفوسفات تقل \omterm{def:b1:plant-water-minerals:mycorrhiza}{الميكوريزا}، وعلى
الغنية بالنترات تصنع البقوليات \omterm{def:b1:plant-water-minerals:nodules}{عقيدات} أقل --- فلا يُبقى على
المقاولين إلا ما داموا أرخص من أداء العمل وحده.
\end{solution}

\begin{solution}{pb:b1:plant-water-minerals:1}
\textbf{1.} $-0.9 + 0.5 = \qty{-0.4}{MPa}$.
\textbf{2.} $-0.2 > -0.4 > -0.6 > -1.3$: كل خطوة أخفض، فينتقل الماء
من التربة $\to$ القشرة $\to$ \omterm{def:b1:flowering-plant-organization:tissues}{الخشب} $\to$ \omterm{def:b1:flowering-plant-organization:organs}{الورقة}.
\textbf{3.} $137\ln 0.5 = \qty{-95}{MPa}$.
\textbf{4.} من التربة إلى \omterm{def:b1:flowering-plant-organization:organs}{الورقة} \qty{1.1}{MPa}؛ ومن \omterm{def:b1:flowering-plant-organization:organs}{الورقة} إلى
الهواء \qty{94}{MPa}: أي \qty{99}{\%} من المجموع هي الخطوة الأخيرة.
\textbf{5.} التربة عند $-0.9$ دون \omterm{def:b1:cell-unit-of-life:cell}{خلايا} القشرة عند $-0.4$: فيغادر
الماء \omterm{def:b1:cell-unit-of-life:cell}{الخلايا} حتى تهبط إلى $-0.9$، أي ينزل امتلاؤها إلى الصفر؛
ويتوقف الالتقاط حتى تخفض \omterm{def:b1:cell-unit-of-life:cell}{الخلايا} $\Psi_s$ فيها بتراكم المنحلات.
\textbf{6.} صارت التربة الآن دون \omterm{def:b1:flowering-plant-organization:organs}{الأوراق} عند $-1.3$: فيجري الماء
خارجًا من النبات؛ وتفقد \omterm{def:b1:flowering-plant-organization:organs}{الأوراق} امتلاءها وتذبل. والسقي يرفع $\Psi$
التربة فوق كمون \omterm{def:b1:flowering-plant-organization:organs}{الجذر} فيستأنف الجريان؛ وتمتلئ \omterm{def:b1:flowering-plant-organization:organs}{الأوراق}، وقد بقيت
جدرها سليمة.
\textbf{7.} $2\times 10^{-7}\times 0.2\times(-0.2 - (-0.6)) =
\qty{1.6e-8}{m^3/s}$.
\textbf{8.} $1.6\times 10^{-8}\times 86\,400 = \qty{1.4e-3}{m^3} =
\qty{1.4}{L}$ في اليوم.
\textbf{9.} \qty{1400}{g} على \qty{0.5}{m^2} وفي \qty{12}{h} من ضوء
النهار (أكثرها): أي نحو \qty{230}{g} في المتر المربع في الساعة.
\textbf{10.} $230/18 = \qty{12.8}{mol}$ في المتر المربع في الساعة، أي \qty{3.5}{mmol} في المتر المربع في الثانية --- أي ألف جزيء ماء مقابل كل $\mathrm{CO_2}$ مثبَّت عند \qty{4}{\micro mol}.
\textbf{11.} الفرق $-0.2 - (-0.4) = 0.2$: فينتصف الدفق، أي
\qty{0.7}{L} في اليوم. وإغلاق الثغور قد قطع الفقد، فارتخى التوتر،
وارتفع السلّم كله: فالنبات يقايض كسب الكربون بالماء حين يشتد جفاف
الهواء.
\textbf{12.} ينتقل الماء نازلًا مع تدرج كمونه هو، وقد أحدثه تبخر
\omterm{def:b1:flowering-plant-organization:organs}{الورقة}، \omterm{def:b1:flowering-plant-organization:organs}{فالجذر} إذن سلبي؛ أما النترات فيجب تركيزها ضد تدرج تركيز
وتدرج كهربائي معًا، ولا يدفع ثمن ذلك إلا الاقتران \omterm{prop:b1:plant-water-minerals:uptake}{بمضخة البروتونات}.
\textbf{13.} $\qty{1.4}{L}\times\qty{2}{mmol/L} = \qty{2.8}{mmol}$، أي
\qty{39}{mg} من النتروجين.
\textbf{14.} نعم، مع فائض قدره الثلث؛ ويُخزن الفائض في الفجوات
نتراتٍ أو يُستبعد عند \omterm{prop:b1:flowering-plant-organization:stemroot}{البشرة الداخلية}.
\textbf{15.} $\Delta G = RT\ln(5/2) + zFV = 2.48\times 0.92 + (-1)(96.5)(-0.15)
= 2.3 + 14.5 = \qty{16.8}{kJ/mol}$.
\textbf{16.} لكل بروتون: $RT\ln(10^{-7.5}/10^{-5.5}) + F(-0.15) =
-11.4 - 14.5 = \qty{-25.9}{kJ}$؛ ولبروتونين: \qty{-52}{kJ}، أي ثلاثة
أمثال المطلوب.
\textbf{17.} $\qty{2.14}{mmol}\times 8 = \qty{17}{mmol}$ من الإلكترونات.
\textbf{18.} $0.5\times 4\times 10^{-6}\times 43\,200 = \qty{86}{mmol}$
من $\mathrm{CO_2}$، أي \qty{346}{mmol} من الإلكترونات: أي \qty{5}{\%}
منها إلى النترات.
\textbf{19.} كان يوفّر الإلكترونات الثمانية لكل نتروجين؛ لكن
التقاط كاتيون وتحرير بروتون مقابل كل واحد منه يحمّض التربة حول
\omterm{def:b1:flowering-plant-organization:organs}{الجذر} وعلى \omterm{def:b1:cell-unit-of-life:cell}{الخلية} أن تصدّر \omterm{def:b1:water-small-molecules:ph}{الحمض}، والأمونيوم سام إن تراكم، فوجب
تمثيله فورًا في \omterm{def:b1:flowering-plant-organization:organs}{الجذر}.
\textbf{20.} \qty{1.07}{mmol} من $\mathrm{N_2}$ $\times 16 =
\qty{17}{mmol}$ من \omterm{def:b1:water-small-molecules:atp}{ATP}.
\textbf{21.} $6\times 0.03 = \qty{0.18}{g}$ من السكر، أي \qty{6}{\%} من
\qty{3}{g}. ومع توافر \qty{2}{mmol/L} من النترات بكلفة اختزالها
(\qty{5}{\%} من الإلكترونات)، يكون \omterm{def:b1:photosynthesis:fixation}{التثبيت} صفقة رديئة وتصنع
البقولية \omterm{def:b1:plant-water-minerals:nodules}{عقيدات} قليلة؛ وبلا نترات يكون هو السبيل الوحيد، ويكون
\qty{6}{\%} من الإنتاج رخيصًا مقابل \omterm{def:b1:proteins:peptide}{بروتين} النبات كله.
\textbf{22.} $1.4\times 2\times 10^{-6} = \qty{2.8}{\micro mol}$، أي
\qty{0.09}{mg} من الفوسفور --- أي جزء من أربعين من الحاجة؛ فعلى
النبات أن يعتمد على الانتشار إلى أشعار جذره، وهي تستنفد ما حولها
سريعًا، وقبل ذلك على فطره الميكوريزي.
\textbf{23.} يسدّ الشريط \omterm{def:b1:plant-water-minerals:pathways}{الأبوبلاست}، فلا تستطيع النترات دخول \omterm{def:b1:flowering-plant-organization:tissues}{الخشب}
إلا عبر \omterm{def:b1:membranes-transport:membrane}{أغشية} \omterm{prop:b1:flowering-plant-organization:stemroot}{البشرة الداخلية} التي تضخها إلى الداخل؛ ويمنع عصير
\omterm{def:b1:flowering-plant-organization:tissues}{الخشب} المركّز من التسرب راجعًا على طول الجدر، فيُحفظ التدرج.
\textbf{24.} كان عصير خشبه يشبه محلول التربة --- ممددًا غير منتقى،
فيه الصوديوم وأيونات أخرى غير مرغوبة --- ولم يكن ليبقى مركّزًا؛ وفي
تربة مالحة كان الصوديوم يبلغ \omterm{def:b1:flowering-plant-organization:organs}{الأوراق} بحرية فيتسمم النبات، حيث
يستبعده النبات السليم عند \omterm{prop:b1:flowering-plant-organization:stemroot}{البشرة الداخلية}.
\textbf{25.} نحو \qty{1.4}{L} من الماء في اليوم عبر \qty{0.2}{m^2}
من \omterm{def:b1:flowering-plant-organization:organs}{الجذر}، تحمل \qty{39}{mg} من النتروجين في مقابل حاجة قدرها
\qty{30}{}، ويأخذ اختزالها نحو \qty{5}{\%} من إلكترونات \omterm{def:b1:flowering-plant-organization:organs}{الورقة}
الضوئية.
\end{solution}
